\begin{abstract}
Continuous Cellular Automata (CCAs) provide an expressive continuous substrate for Artificial Life, yet computational simulations of Open-Ended Evolution (OEE) routinely encounter premature convergence, severe parameter fragility, and the simulation gap. Flow Lenia addresses numerical mass dissipation by coupling continuous cellular dynamics with mass-conserving fluid advection partial differential equations. However, unguided continuous parameter spaces remain over 99.9\% sterile, while uniform environments lack the selective friction necessary to drive functional morphological adaptation.

This thesis presents a GPU-accelerated research framework in native JAX that systematically evaluates the impact of structured spatial heterogeneity on phenotypic adaptation, soft biomechanics, and automated discovery in Flow Lenia. The simulation engine combines two-dimensional spectral Fourier convolutions, multi-shell concentric ring kernels, and semi-Lagrangian bilinear reintegration tracking \parencite{moroz2020reintegration, plantec2025flowlenia}; across our verification benchmarks, this formulation preserved global mass within floating-point precision with zero artificial coordinate damping. To navigate the continuous parameter space without human template bias, we deploy an Intrinsically Motivated Goal Exploration Process (IMGEP) across a three-dimensional metric space spanning center-of-mass motility, evolutionary activity, and compression complexity, evaluated through a sequential five-stage morphological quality filter.

Empirical evaluations across structured environmental assays indicate that spatial heterogeneity fosters structural plasticity and destabilizes trivial static equilibria in the evaluated regimes:
\begin{enumerate}
    \item \textbf{Chemotactic Phase Transitions:} Directional nutrient gradients expose a cohesion-fission transition governed by internal surface tension and external tensile pull, distinguishing unitary elastomeric migration from shear-induced amoeboid mitosis.
    \item \textbf{Elastomeric Aperture Transmission:} A soft-bodied barrier constriction assay measures the aperture transmission coefficient $T(W)$, showing a sigmoidal transmission profile from complete obstruction ($T(8\text{ px}) = 0.0\%$) to high-aperture chamber traversal ($T(64\text{ px}) = 77.0\%$) with an empirical necking interval near $W \approx 28\text{--}32\text{ px}$ in our sample ($n=3$).
    \item \textbf{Dynamic Niche Construction:} Localized resource depletion creates a negative feedback loop that eliminates static crystal attractors in our evaluated sample ($0.0\%$ resting equilibria), compelling continuous cyclic foraging trajectories.
    \item \textbf{Macro-Scale Ecology:} Across 22,500 continuous simulation steps in a multi-chamber arena benchmark, categorical Gumbel-Max flux sampling mitigates parameter homogenization and sustains distinct territorial lineages.
\end{enumerate}

Furthermore, we examine two practical failure modes in automated discovery: \textit{stochastic seed sensitivity} (where transient initialization gradients induce false-positive motility) and the \textit{solidity-motility trade-off} (where strict solidity thresholds favor static crystals over dynamic gliders). Finally, we establish a clear epistemological demarcation: in spatial navigation tasks, topological pathfinding is solved by external potential fields while Flow Lenia acts as the deformable soft-matter physical substrate. We conclude that while continuous cellular automata do not exhibit autonomous cognitive agency, coupling mass-conserving fluid mechanics with structured spatial friction provides a grounded, scalable framework for modeling soft-bodied dynamics, emergent morphology, and decentralized artificial life.
\end{abstract}